\section{RMSProp}

\subsection{Fundamentos teóricos}

RMSProp adapta el paso a la escala reciente de los gradientes de cada
coordenada \cite{tieleman2012rmsprop,d2l_rmsprop}. En lugar de aplicar la
misma magnitud de ajuste a todas las componentes, mantiene un promedio
móvil exponencial de los gradientes al cuadrado:
\[
s_{t+1} = \gamma s_t + (1-\gamma)g_t\odot g_t, \qquad
x_{t+1} = x_t - \frac{\alpha g_t}{\sqrt{s_{t+1}+\epsilon}},
\]
con $g_t=\nabla f(x_t)$, $0\leq\gamma<1$ y operaciones elemento a elemento.
Un gradiente reciente grande aumenta el denominador de su coordenada y
reduce el paso correspondiente; $\epsilon>0$ evita la división entre cero.
El promedio móvil conserva información reciente sin acumular
indefinidamente todos los gradientes anteriores.

\subsection{Formulación de la solución}

La función
\texttt{run\_rmsprop(alpha, gamma, epsilon, t, func, initial\_point)}
inicializa $s_0=0$, calcula el gradiente con Autograd igual que GD y realiza
ambas actualizaciones como operaciones tensoriales sobre las dos
coordenadas a la vez. Si el valor, el gradiente o el promedio móvil dejan
de ser finitos, lanza un \texttt{ValueError}. En todo el trabajo se usa
$\epsilon=10^{-6}$.

\subsection{Demostración con hiperparámetros fijados manualmente}

Se usó $\gamma=0.9$ y $\epsilon=10^{-6}$ en las tres funciones, con $P=25$
iteraciones y el mismo punto inicial que la demostración de GD
(Tabla~\ref{tab:gd_demo}):

\begin{table}[H]
\centering
\begin{tabular}{lrrr}
\toprule
Función & $\alpha$ & $f(x_0)$ & $f(x_{25})$ \\
\midrule
$f_0$ & $0.5$  & $58.030396$   & $0.000233$ \\
$f_1$ & $0.05$ & $14.279894$   & $12.579832$ \\
$f_2$ & $0.4$  & $2419.764335$ & $0.036919$ \\
\bottomrule
\end{tabular}
\caption{Demostración de RMSProp con 25 iteraciones, desde los mismos
puntos iniciales que GD.}
\label{tab:rms_demo}
\end{table}

En $f_1$ el valor final permanece lejos de cero: el punto inicial condujo,
igual que en GD, a la cuenca del mínimo local cercano a $(0,\ 7)$. En $f_2$,
RMSProp llega cerca del mismo mínimo que GD, $(-3.779310,\ -3.283186)$, con
$\alpha=0.4$, un paso 200 veces mayor que el de GD, sin divergir: al
dividir entre $\sqrt{s_t}$, la magnitud de cada componente del paso queda
acotada aproximadamente por $\alpha/\sqrt{1-\gamma}$ aunque el gradiente
sea del orden de miles. La Figura~\ref{fig:rms_demo} muestra estas
corridas.

\begin{figure}[H]
\centering
\begin{subfigure}{0.49\linewidth}
\includegraphics[width=\linewidth]{curva-de-aprendizaje-de-rmsprop-sobre-f0}
\caption{Curva de aprendizaje en $f_0$.}
\end{subfigure}
\hfill
\begin{subfigure}{0.36\linewidth}
\includegraphics[width=\linewidth]{trayectoria-de-rmsprop-sobre-f0}
\caption{Trayectoria en $f_0$.}
\end{subfigure}

\begin{subfigure}{0.49\linewidth}
\includegraphics[width=\linewidth]{curva-de-aprendizaje-de-rmsprop-sobre-f1}
\caption{Curva de aprendizaje en $f_1$.}
\end{subfigure}
\hfill
\begin{subfigure}{0.36\linewidth}
\includegraphics[width=\linewidth]{trayectoria-de-rmsprop-sobre-f1}
\caption{Trayectoria en $f_1$.}
\end{subfigure}

\begin{subfigure}{0.49\linewidth}
\includegraphics[width=\linewidth]{curva-de-aprendizaje-de-rmsprop-sobre-f2}
\caption{Curva de aprendizaje en $f_2$.}
\end{subfigure}
\hfill
\begin{subfigure}{0.36\linewidth}
\includegraphics[width=\linewidth]{trayectoria-de-rmsprop-sobre-f2}
\caption{Trayectoria en $f_2$.}
\end{subfigure}
\caption{Curvas de aprendizaje y trayectorias de la demostración de
RMSProp (Tabla~\ref{tab:rms_demo}).}
\label{fig:rms_demo}
\end{figure}

\subsection{Calibración de hiperparámetros con Optuna}

Optuna ajusta $\alpha$ (escala logarítmica) y $\gamma\in[0.5,\ 0.999]$;
$\epsilon$ queda fijo en $10^{-6}$ y las iteraciones en 25. Los rangos de
$\alpha$ fueron $[0.01,\ 1]$ para $f_0$, $[0.001,\ 0.5]$ para $f_1$ y
$[0.01,\ 2]$ para $f_2$. El protocolo es el mismo que en GD: 25 ensayos,
6 puntos de ajuste y el promedio de $f$ final como objetivo. Ningún ensayo
de RMSProp se podó.

\begin{table}[H]
\centering
\begin{tabular}{lrrcrr}
\toprule
Función & Mejor $\alpha$ & Mejor $\gamma$ & Convergencias & Iter.\ promedio & Promedio $f$ final \\
\midrule
$f_0$ & $0.629157$ & $0.992523$ & 6/6 & $5.00$  & $8.31\times10^{-19}$ \\
$f_1$ & $0.018006$ & $0.998647$ & 0/6 & ---     & $13.035986$ \\
$f_2$ & $0.341559$ & $0.611311$ & 2/6 & $11.50$ & $2.656585$ \\
\bottomrule
\end{tabular}
\caption{Mejores hiperparámetros de RMSProp encontrados por Optuna (25
ensayos, 25 iteraciones, 6 puntos de ajuste).}
\label{tab:rms_optuna}
\end{table}

En $f_1$, Optuna eligió un paso pequeño y $\gamma$ cercano a 1: con esa
configuración RMSProp desciende de forma estable hasta el mínimo local más
cercano. Ninguna combinación probada alcanzó la tolerancia en Ackley, y el
promedio obtenido ($13.04$) es peor que el de GD ($6.27$). La
Figura~\ref{fig:rms_optuna} muestra el proceso de calibración.

\begin{figure}[H]
\centering
\begin{subfigure}{0.36\linewidth}
\includegraphics[width=\linewidth]{historial-de-optimizacion-rmsprop-sobre-f0}
\caption{Historial de optimización en $f_0$.}
\end{subfigure}
\hfill
\begin{subfigure}{0.62\linewidth}
\includegraphics[width=\linewidth]{efecto-de-hiperparametros-rmsprop-sobre-f0}
\caption{Efecto de $\alpha$ y $\gamma$ en $f_0$.}
\end{subfigure}

\begin{subfigure}{0.36\linewidth}
\includegraphics[width=\linewidth]{historial-de-optimizacion-rmsprop-sobre-f1}
\caption{Historial de optimización en $f_1$.}
\end{subfigure}
\hfill
\begin{subfigure}{0.62\linewidth}
\includegraphics[width=\linewidth]{efecto-de-hiperparametros-rmsprop-sobre-f1}
\caption{Efecto de $\alpha$ y $\gamma$ en $f_1$.}
\end{subfigure}

\begin{subfigure}{0.36\linewidth}
\includegraphics[width=\linewidth]{historial-de-optimizacion-rmsprop-sobre-f2}
\caption{Historial de optimización en $f_2$.}
\end{subfigure}
\hfill
\begin{subfigure}{0.62\linewidth}
\includegraphics[width=\linewidth]{efecto-de-hiperparametros-rmsprop-sobre-f2}
\caption{Efecto de $\alpha$ y $\gamma$ en $f_2$.}
\end{subfigure}
\caption{Proceso de calibración de RMSProp con Optuna para $f_0$, $f_1$ y
$f_2$ (de arriba hacia abajo): historial de optimización (izquierda) y
efecto de $\alpha$ y $\gamma$ con el mejor ensayo marcado (derecha).}
\label{fig:rms_optuna}
\end{figure}

Como evidencia de los mejores valores (punto 2a), se repitieron las 6
corridas de ajuste con los valores de la Tabla~\ref{tab:rms_optuna}
(Figura~\ref{fig:rms_curvas_calibracion}). La curva negra es el promedio,
cuyo último valor coincide con el objetivo de Optuna:
$8.31\times10^{-19}$, $13.035986$ y $2.656585$. En $f_0$ las seis corridas
convergen, en 5 iteraciones en promedio. En $f_1$ cada corrida desciende de
forma estable hasta un mínimo local cercano a su inicio, con valores entre
$10$ y $16$. En $f_2$ el paso acotado hace que el descenso sea más gradual
que en GD; todas las corridas terminan con $f$ entre $1.2$ y $4.7$,
incluidas las dos que alcanzaron la tolerancia (en las iteraciones 11 y 12)
y después se alejaron del mínimo por las oscilaciones del paso normalizado.

\begin{figure}[H]
\centering
\begin{subfigure}{0.70\linewidth}
\includegraphics[width=\linewidth]{curvas-de-aprendizaje-de-rmsprop-con-los-mejores-hiperparametros-sobre-f0}
\caption{$f_0$, $\alpha=0.629157$, $\gamma=0.992523$.}
\end{subfigure}

\begin{subfigure}{0.70\linewidth}
\includegraphics[width=\linewidth]{curvas-de-aprendizaje-de-rmsprop-con-los-mejores-hiperparametros-sobre-f1}
\caption{$f_1$, $\alpha=0.018006$, $\gamma=0.998647$.}
\end{subfigure}

\begin{subfigure}{0.70\linewidth}
\includegraphics[width=\linewidth]{curvas-de-aprendizaje-de-rmsprop-con-los-mejores-hiperparametros-sobre-f2}
\caption{$f_2$, $\alpha=0.341559$, $\gamma=0.611311$.}
\end{subfigure}
\caption{Curvas de aprendizaje de RMSProp con los mejores
hiperparámetros sobre los 6 puntos de ajuste, para $f_0$, $f_1$ y $f_2$.}
\label{fig:rms_curvas_calibracion}
\end{figure}

\subsection{Evaluación con los mejores hiperparámetros}

Con los valores de la Tabla~\ref{tab:rms_optuna}, desde los mismos 10
puntos y con un máximo de 50 iteraciones (Sección~\ref{sec:evaluacion}):

\begin{itemize}
    \item $f_0$: 10/10 corridas convergen, con $4.20$ iteraciones en
    promedio; promedio de $f$ minimizada $3.67\times10^{-4}$.
    \item $f_1$: 0/10 corridas convergen; promedio de $f$ final $12.4341$.
    Los valores finales, entre $8.6165$ y $15.5491$, corresponden a mínimos
    locales alejados del origen: con el paso pequeño calibrado, RMSProp
    desciende hasta un mínimo local cercano a su punto inicial, como muestra
    la trayectoria de la corrida 5 (Figura~\ref{fig:rms_mejores}).
    \item $f_2$: 5/10 corridas convergen, con $14.20$ iteraciones en
    promedio; promedio de $f$ minimizada $1.2026$.
\end{itemize}

Las 5 corridas de $f_2$ que no convergen terminan con valores entre
$1.9068$ y $3.8037$. Con $\gamma=0.611311$ el promedio $s_t$ sigue de cerca
a $g_t^2$, de modo que el paso normalizado no se reduce cuando el gradiente
se hace pequeño y la trayectoria oscila alrededor de un mínimo. La
corrida 5 ilustra este comportamiento (Tabla~\ref{tab:silla}): llega a
$f=0.0170$ en $t=10$ y después vuelve a valores cercanos a $2.25$.

\subsection{Mejor corrida por función}

\begin{figure}[H]
\centering
\begin{subfigure}{0.36\linewidth}
\includegraphics[width=\linewidth]{evaluacion-mejor-trayectoria-de-rmsprop-sobre-f0}
\caption{Trayectoria en $f_0$ (corrida 7).}
\end{subfigure}
\hfill
\begin{subfigure}{0.49\linewidth}
\includegraphics[width=\linewidth]{evaluacion-curva-de-aprendizaje-de-rmsprop-sobre-f0}
\caption{Curva de aprendizaje en $f_0$.}
\end{subfigure}

\begin{subfigure}{0.36\linewidth}
\includegraphics[width=\linewidth]{evaluacion-mejor-trayectoria-de-rmsprop-sobre-f1}
\caption{Trayectoria en $f_1$ (corrida 5).}
\end{subfigure}
\hfill
\begin{subfigure}{0.49\linewidth}
\includegraphics[width=\linewidth]{evaluacion-curva-de-aprendizaje-de-rmsprop-sobre-f1}
\caption{Curva de aprendizaje en $f_1$.}
\end{subfigure}

\begin{subfigure}{0.36\linewidth}
\includegraphics[width=\linewidth]{evaluacion-mejor-trayectoria-de-rmsprop-sobre-f2}
\caption{Trayectoria en $f_2$ (corrida 5).}
\end{subfigure}
\hfill
\begin{subfigure}{0.49\linewidth}
\includegraphics[width=\linewidth]{evaluacion-curva-de-aprendizaje-de-rmsprop-sobre-f2}
\caption{Curva de aprendizaje en $f_2$.}
\end{subfigure}
\caption{Mejores corridas de RMSProp. $f_0$: corrida 7, converge en 2
iteraciones ($f=4.43\times10^{-5}$). $f_1$: corrida 5, no converge y queda
en el mínimo local cercano a $(4,\ 0)$ ($f=8.616490$). $f_2$: corrida 5,
converge en 7 iteraciones ($f=8.29\times10^{-4}$).}
\label{fig:rms_mejores}
\end{figure}

\subsection{¿Por qué RMSProp evita mejor los puntos silla que GD?}
\label{sec:silla}

En un punto silla el gradiente es cero, y en su vecindad es pequeño en
todas las direcciones. El paso de GD es proporcional al gradiente,
$\alpha\nabla f$, así que cerca de un punto silla los pasos se vuelven
diminutos y el algoritmo puede pasar muchas iteraciones en una meseta,
aunque exista una dirección de curvatura negativa por la que podría
descender \cite{dauphin2014identifying,goodfellow2016deep}.

RMSProp divide cada componente del gradiente entre la raíz de su promedio
reciente de cuadrados. Si el gradiente se reduce de forma sostenida, $s_t$
también se reduce y el cociente $g_t/\sqrt{s_t+\epsilon}$ se mantiene del
orden de 1: el tamaño del paso depende principalmente de $\alpha$ y no de la
magnitud del gradiente. Además, la normalización es por coordenada, por lo
que una dirección con gradiente pequeño (por ejemplo, la dirección de
escape de la silla) recibe un paso relativamente mayor que en GD
\cite{tieleman2012rmsprop,ruder2016overview}. En consecuencia, RMSProp
atraviesa las regiones planas alrededor de los puntos silla en menos
iteraciones.

La evaluación contiene un caso que lo ilustra. La corrida 5 inicia en
$(4.4551,\ 0.2191)$, cerca del punto silla $(3.385154,\ 0.073852)$ de $f_2$,
cuyo valor es $13.311926$:

\begin{table}[H]
\centering
\small
\begin{tabular}{lrrrrrrrrr}
\toprule
Algoritmo & Dist.\ mín. & $t$ & \multicolumn{7}{c}{$f_2(x_t)$} \\
\cmidrule(lr){4-10}
 & & & $t=0$ & $t=5$ & $t=10$ & $t=20$ & $t=30$ & $t=40$ & $t=50$ \\
\midrule
GD      & $0.043$ & 5 & $88.4406$ & $13.3026$ & $13.2850$ & $13.1924$ & $12.7921$ & $11.2136$ & $6.7753$ \\
RMSProp & $0.659$ & 1 & $88.4406$ & $0.3466$  & $0.0170$  & $2.9322$  & $2.2754$  & $2.2485$  & $2.2541$ \\
\bottomrule
\end{tabular}
\caption{Corrida 5 sobre $f_2$ con los mejores hiperparámetros: menor
distancia de la trayectoria al punto silla, iteración en que ocurre y valor
de $f_2$ en varias iteraciones.}
\label{tab:silla}
\end{table}

GD llega a $0.043$ unidades del punto silla en $t=5$ y permanece con
$f_2\approx13.3$, prácticamente el valor de la silla, hasta cerca de
$t=20$; en 50 iteraciones no alcanza a converger. RMSProp, desde el mismo
punto, pasa lejos de la silla y converge en 7 iteraciones. Es un solo caso
y no un experimento controlado, pero es consistente con la explicación
anterior.
